% Chapter 8: Entorno Regulatorio y de Políticas
\chapter{Entorno Regulatorio y de Políticas}

\subsection{Regulaciones Clave por Geografía}

El cumplimiento ambiental y de obtención de permisos (permitting) es el principal factor de retraso en proyectos mineros globales. El marco regulatorio en las geografías objetivo se detalla a continuación:

\begin{enumerate}
    \item \textbf{Chile:} Regulado principalmente por el Servicio de Evaluación Ambiental (SEA). Las exigencias respecto a la huella hídrica y de carbono son sumamente estrictas. La nueva Ley de Royalty Minero (vigente) reestructura el impuesto a los explotadores mineros, elevando la carga tributaria efectiva pero dando certeza jurídica de largo plazo.
    
    \item \textbf{Perú:} El Senace (Servicio Nacional de Certificación Ambiental para las Inversiones Sostenibles) es el encargado de evaluar los EIA detallados. Las regulaciones de relacionamiento comunitario e consulta previa son críticas debido a la alta conflictividad social en el corredor minero del sur.
    
    \item \textbf{Brasil:} La ANM (Agência Nacional de Mineração) y agencias ambientales estatales (como SEMAD en Minas Gerais) regulan la actividad. Tras los incidentes históricos de relaves, la regulación de presas de relaves upstream está estrictamente prohibida, obligando a las empresas a rediseñar sus plantas para disposición de relaves secos.
    
    \item \textbf{Estados Unidos:} El marco federal (NEPA) coordinado por el Bureau of Land Management (BLM) e EPA genera tiempos de obtención de permisos de entre 7 y 10 años. Sin embargo, bajo la administración actual, existen iniciativas y presiones ejecutivas para acelerar la tramitación ambiental de proyectos de litio y cobre mediante ``fast-tracking'' por considerarse de seguridad nacional.
    
    \item \textbf{Rusia:} Sujeto a regulaciones locales rusas y fuertemente impactado por sanciones occidentales externas que limitan la transferencia tecnológica y la exportación de equipos mineros avanzados.
\end{enumerate}

\subsection{Cumplimiento y Oportunidad en Permitting}

Las consultoras con profundo conocimiento de los marcos regulatorios locales y con capacidades de tramitación de permisos ambientales (permitting) disfrutan de una demanda inelástica. Las mineras están dispuestas a pagar tarifas premium para acortar en 6 o 12 meses el ciclo de aprobación ambiental de un proyecto multimillonario.
